%%%%%%%%%%%%%%%%%%%%%%%%%%%%%%%%
\chapter{フィッツヒュー・南雲方程式に特化した手法}\label{Proposing}
%%%%%%%%%%%%%%%%%%%%%%%%%%%%%%%%

%%%%%%%%%%%%%%%%%%%%%%%%%%%%%%%%%%%%%%%%%%%%%%%%%%%%%%%%%%
                        %第四章%     %フィッツヒュー・南雲方程式に特化した手法%
%%%%%%%%%%%%%%%%%%%%%%%%%%%%%%%%%%%%%%%%%%%%%%%%%%%%%%%%%%

\section{\(\tilde{\mathcal{F}}\)ノルム}
%%%%%%%%%%%%%%%%% [[[[Fチルダノルム]]]] %%%%%%%%%%%%%%%%%%%%%%%%%%%%%%%%
\(\qty(u,v)_{H^1_0}\coloneqq\qty(\nabla u,\nabla v)\)を用いて，不安定解をもつ方程式に適用できる\(\tilde{\mathcal{F}}\)内積を導入する．
この時，\(H^1_0\qty(\Omega)\)上の内積を次のように定める．
\begin{align}\label{eq:4章f.tilde}
    \qty(u,v)_{\tilde{\mathcal{F}}}\coloneqq\qty(\nabla u,\nabla v)_{L^2}+\qty(\dfrac{1}{\epsilon^2}(\alpha + 3\hat{u}^2\gamma +\delta \mathcal{B})u,v)_{L^2}
\end{align}
よって，\eqref{eq:4章f.tilde}が内積の公理を満たせばよい．
\eqref{eq:4章f.tilde}の場合，
\(\hat{u}^2 \geq 0,\alpha \geq 0,\gamma \geq 0, \delta \mathcal{B} \geq 0\)
であるため，右辺が\((\nabla u,\nabla v)_{L^2}+(\frac{1}{\epsilon^2}(\alpha + 3\hat{u}^2\gamma +\delta \mathcal{B})u,v)_{L^2}\geq 0\)となる事から，内積の公理を満たすことがわかる．
この内積から誘導されるノルムを
\begin{align}
    \norm{u}_{\tilde{\mathcal{F}}} \coloneqq \sqrt{(u,u)_{\tilde{\mathcal{F}}}}
\end{align}
で定義する.
ここで，\(W\)を\(H^1_0(\Omega)\)に\(\tilde{F}\)の内積を導入したHilbert空間とする．
また，\(W_N\)を\(W\)の有限次元部分空間とし，\(W_\perp\)を\(W\)の\(W_N\)に対する直交補空間とする．
また，\(W^{-1}\)を\(W\)の共役空間とし，そのノルムを
\begin{align}
  \norm{u}_{\tilde{\mathcal{F}}^{-1}}
    =
   \sup\limits_{\phi \in H^1_0(\Omega) \backslash \{ 0 \} }
   \frac{|\langle u,v \rangle|}{\norm{v}_{\tilde{\mathcal{F}}}}
\end{align}
とする．
そのノルムで Newton-Kantorovich の定理を考える．

%\begin{align}
%    \sup\limits_{\phi \in H^1_0(\Omega) \backslash \{ 0 \} }
%    \frac{|( T,v )|}{\norm{v}_{\tilde{\mathcal{F}}}}
%\end{align}
\section{\(K\)の評価}%えーたの評価までここで述べる% %C_{\tilde{\mathcal{F}}}ものべる%
%%%%%%%%%%%%%%%%% [[[[Kの評価]]]] %%%%%%%%%%%%%%%%%%%%%%%%%%%%%%%%
\(\norm{\mathcal{F}'\qty[\hat{u}]^{-1}}_{L(W^{-1},W)} \leq K\)
を満たす\(K\)を評価する．\\
線形作用素\(\mathcal{A}_{\tilde{\mathcal{F}}}\coloneqq H^1_0(\Omega)\rightarrow H^{-1}(\Omega)\)を
\begin{align}
   \langle \mathcal{A}_{\tilde{\mathcal{F}}} u , w \rangle 
   \coloneqq 
   \qty(u,v)_{\tilde{\mathcal{F}}}
  %  (\nabla u , \nabla w)_{L^2}
  %  +\qty(\dfrac{1}{\epsilon^2}(\alpha + 3\hat{u}^2\gamma +\delta \mathcal{B})u , w)_{L^2}
\end{align}
とすると，作用素ノルムの定義より，\(\norm{\mathcal{F}'\qty[\hat{u}]^{-1}}_{L\qty(W^{-1},W)} \)
は次のように表すことができる．
\begin{align}
    \norm{\mathcal{F}'\qty[\hat{u}]^{-1}}_{L\qty(W^{-1},W)} 
    &=
    \sup\limits_{\phi \in H^{-1}(\Omega) \backslash \{ 0 \} }
    \dfrac{\norm{\mathcal{F}'\qty[\hat{u}]^{-1}\phi}_{\tilde{\mathcal{F}}}}{\norm{\phi}_{\tilde{\mathcal{F}}^{-1}}}\\
    &=
    \sup\limits_{\phi \in H^1_0 (\Omega)\backslash \{ 0 \} }
    \dfrac{\norm{\phi}_{\tilde{\mathcal{F}}}}{\norm{\mathcal{F}'\qty[\hat{u}]\phi}_{\tilde{\mathcal{F}}^{-1}}}\\
    &=
    \sup\limits_{\phi \in H^1_0 (\Omega)\backslash \{ 0 \} }
    \dfrac{\norm{\phi}_{\tilde{\mathcal{F}}}}{\norm{\mathcal{A}_{\tilde{\mathcal{F}}}^{-1} \mathcal{F}'\qty[\hat{u}] \phi}_{\tilde{\mathcal{F}}}}\\
\end{align}
ここで，\(\mathcal{A}_{\tilde{\mathcal{F}}}^{-1} \mathcal{F}'\qty[\hat{u}]\)についてスペクトル分解すると
\(\mathcal{A}_{\tilde{\mathcal{F}}}^{-1} \mathcal{F}'\qty[\hat{u}] = \int \mu d E\)
と置き換えられる．\\
したがって，
\begin{align}
  \norm{\mathcal{F}'\qty[\hat{u}]^{-1}}_{L\qty(W^{-1},W)} 
  &=
  \sup\limits_{\phi \in H^1_0 (\Omega)\backslash \{ 0 \} }
  \dfrac{\norm{\phi}_{\tilde{\mathcal{F}}}}{\norm{\int \mu d E\phi}_{\tilde{\mathcal{F}}}}\\
  &\leq
  \sup\limits_{\phi \in H^1_0 (\Omega)\backslash \{ 0 \} }
  \dfrac{\norm{\phi}_{\tilde{\mathcal{F}}}}{|\mu |\norm{\int d E\phi}_{\tilde{\mathcal{F}}}}\\
  &\leq
  \dfrac{1}{|\mu_1|}
\end{align}
と評価できる．このとき，\(\mu\) はスカラとして取り出すことができ， 
絶対値をつけることで右辺が大きくなる．
なお，\(\mu_1\) はゼロに最も近い絶対最小固有値を意味する．
ここで，非線形作用素\(\mathcal{N}_{\tilde{\mathcal{F}}}\coloneqq H^1_0(\Omega)\rightarrow H^{-1}(\Omega)\)を
\begin{align}
  \langle \mathcal{N}_{\tilde{\mathcal{F}}} (u) , w \rangle
  \coloneqq \qty(\frac{f(u)-\delta \mathcal{B}u}{\epsilon^2} , w)_{L^2}
  +\qty(\dfrac{1}{\epsilon^2}(\alpha + 3\hat{u}^2\gamma +\delta \mathcal{B})u , w)_{L^2},\;\;\forall w \in H^1_0(\Omega)
\end{align}
とする．
固有値 \(\mu\) の下界評価にあたり，
\begin{align}
  \qty(\mathcal{A}_{\tilde{\mathcal{F}}} - \mathcal{N}'_{\tilde{\mathcal{F}}} [\hat{u}])u =\mu \mathcal{A}_{\tilde{\mathcal{F}}} u
\end{align}
を考える．
この時
\begin{align}
  \dfrac{1}{\mu} = \dfrac{\eta}{\eta - 1}
\end{align}
とすることで，
\begin{align}\label{eq:4固有値前}
  \mathcal{A}_{\tilde{\mathcal{F}}} u = \eta \mathcal{N'}_{\tilde{\mathcal{F}}}[\hat{u}] u 
\end{align}
とすることができる．また，新しく\(d\)内積を以下のように定義する．
\begin{align}
   (u,w)_d \coloneqq \qty(\dfrac{f'[u]-\delta \mathcal{B}}{\epsilon^2}u,w)_{L^2}+\qty(\dfrac{1}{\epsilon^2}(\alpha + 3\hat{u}^2\gamma +\delta \mathcal{B})u,w)_{L^2}
\end{align}
\eqref{eq:4固有値前}を固有値問題に書き直すと
\begin{align}
  \text{Find} \;\;\;\eta \in \mathbb{R}\;\;\; \text{and} \;u\in H^1_0(\Omega)\;\text{s.t.}\;(u, w)_{\tilde{\mathcal{F}}} = \eta (u,w)_d,\;\forall w \in H^1_0 (\Omega)
\end{align}
のスペクトル\(\eta\)は固有値となり，\(\qty{\eta_i}^\infty_{i=1}\)とする．
また離散化すると
\begin{align}
  \text{Find} \;\;\;\eta^h \in \mathbb{R} \;\;\text{and} \;\;u_h\in W_h\;\text{s.t.}\;(u_h, w_h)_{\tilde{\mathcal{F}}} = \eta^h (u_h,w_h)_d,\;\forall w_h \in W_h (\Omega)
\end{align}
となり，その固有値を\(\qty{\eta^h_i}^n_{i=1}\)とする．\\
直交射影\(\mathcal{P}_{h,\tilde{\mathcal{F}}}:H^1_0(\Omega)\rightarrow W_h\)は
\(\qty(u-\mathcal{P}_{h,\tilde{\mathcal{F}} }u , \phi_h)_{\tilde{\mathcal{F}}} = 0,\phi_h\in W_N\)
を満たすとする．定数\(C_{h,d}\)は不等式
\begin{align}\label{eq:C&v-p}
  \norm{v-\mathcal{P}_{h,\tilde{\mathcal{F}}}v}_d 
  \leq C_{h,d} 
  \norm{v-\mathcal{P}_{h,\tilde{\mathcal{F}}}v}_{\tilde{\mathcal{F}}},
  \;\;\forall x \in H^1_0(\Omega)
\end{align}
を満たすとする．そのとき，劉・大石の定理より，\(i \leq n\)となる整数\(i\)番目の固有値\(\eta_i\)は
\begin{align}
  \dfrac{\eta^h_i}{1+C_{h,d}^2\eta^h_i} \leq \eta_i \leq  \eta^h_i
\end{align}
を満たす．

\(C_{h,d}\)を求めるために，\eqref{eq:C&v-p}の両辺にある，\\
\(
  \norm{v-\mathcal{P}_{h,\tilde{\mathcal{F}}}v}_d ,\;\;\; \norm{v-\mathcal{P}_{h,\tilde{\mathcal{F}}}v}_{\tilde{\mathcal{F}}}
\)
の関係式を作る必要がある．

任意の\( v \in H^1_0(\Omega)\)について，
\begin{align}
  {\norm{v}_d}^2 =(v,v)_d =(f'[\hat{u}]v,v)_{L^2}+\qty(\dfrac{1}{\epsilon^2}(\alpha + 3\hat{u}^2\gamma +\delta \mathcal{B})v,v)_{L^2}
\end{align}
である．
このとき\(f'[\hat{u}]\)は，
\begin{align}
  f'[\hat{u}]=&\dfrac{1}{\epsilon^2}\qty(-\alpha+2\beta\hat{u}-3\gamma\hat{u}^2-\delta\mathcal{B})\\
\end{align}
であるため，
\begin{align}
  {\norm{v}_d}^2 &=(v,v)_d =(f'[\hat{u}]v,v)_{L^2}+\qty(\dfrac{1}{\epsilon^2}(\alpha + 3\hat{u}^2\gamma +\delta \mathcal{B})v,v)_{L^2}\\
  &=((\dfrac{1}{\epsilon^2}\qty(2\beta\hat{u}))v,v)_{L^2}\\
  &=\dfrac{1}{\epsilon^2}\norm{2\beta\hat{u}}_{L^\infty}\norm{v}_{L^2}^2
\end{align}
となる．
よって，
\begin{align}
  \norm{v-\mathcal{P}_{h,\tilde{\mathcal{F}}}v}_d 
  \leq&\sqrt{\dfrac{1}{\epsilon^2}\norm{2\beta\hat{u}}_{L^\infty}}\norm{v-\mathcal{P}_{h,\tilde{\mathcal{F}}}v}_{L^2}\\
\end{align}
さらに
\begin{align}
  \norm{v-\mathcal{P}_{h,\tilde{\mathcal{F}}}v}_{L^2}
  \leq
  C_{h,\tilde{\mathcal{F}}}\norm{v-\mathcal{P}_{h,\tilde{\mathcal{F}}}v}_{\tilde{\mathcal{F}}}
\end{align}
が成り立つので
\begin{align}
  \norm{v-\mathcal{P}_{h,\tilde{\mathcal{F}}}v}_d\leq C_{h,\tilde{\mathcal{F}}}\sqrt{\dfrac{1}{\epsilon^2}\norm{2\beta\hat{u}}_{L^\infty}}\norm{v-\mathcal{P}_{h,\tilde{\mathcal{F}}}v}_{\tilde{\mathcal{F}}}
\end{align}
となる．よって
\begin{align}
  C_{h,d}\coloneqq C_{h,\tilde{\mathcal{F}}}\sqrt{\dfrac{1}{\epsilon^2}\norm{2\beta\hat{u}}_{L^\infty}}
\end{align}
と評価できる．

\section{\(C_{h,\tilde{\mathcal{F}}}\)の評価}
\(L\)を\(L_{i,j}\coloneqq(\phi_i,\phi_j)_{L^2}\)，
\(G\)を
\(G\coloneqq (\nabla\phi_i,\nabla\phi_j)_{L^2}
+\qty(\dfrac{1}{\epsilon^2}(\alpha+3\gamma \hat{u}^2 + \delta \mathcal{B} )\phi_i,\phi_j)_{L^2}\)
とすると，\(C_{h,\tilde{\mathcal{F}}}\)は
\begin{align}
  C_{h,\tilde{\mathcal{F}}} \coloneqq 
  C_N\qty(1+\norm{L^{\frac{1}{2}}G^{-1}L^{\frac{T}{2}}}_E
  \norm{\dfrac{1}{\epsilon^2}(\alpha+3\gamma \hat{u}^2 + \delta \mathcal{B} )}_{L(L^2)})
\end{align}
となる．

\section{\(C_{p\tilde{\mathcal{F}}}\)の評価}
定数\(C_{p\tilde{\mathcal{F}}}\)は
\begin{align}
   \norm{u}_{L^2}\leq C_{p\tilde{\mathcal{F}}} \norm{u}_{\tilde{\mathcal{F}}}
\end{align}
で定義される．これは
\begin{align}
  C_{p\tilde{\mathcal{F}}} \coloneqq 
  \sqrt{\dfrac{1+C_{h,\tilde{\mathcal{F}}}^2\norm{L^{\frac{1}{2}}G^{-1}L^{\frac{T}{2}}}_E}{\norm{L^{\frac{1}{2}}G^{-1}L^{\frac{T}{2}}}_E}}
\end{align}
で評価できる．
\section{\(R\)の評価}
%%%%%%%%%%%%%%%%% [[[[Rの評価]]]] %%%%%%%%%%%%%%%%%%%%%%%%%%%%%%%%
\(\norm{\mathcal{F}'(\hat{u})}_{W^{-1}}\leq \delta\)を評価する．
\(\norm{\mathcal{F}'(\hat{u})}_{W^{-1}}\)はノルム\(\norm{\cdot}_{\tilde{\mathcal{F}}^{-1}}\)の定義より以下のように変形できる．

\begin{align}
    \norm{\mathcal{F}\qty(\hat{u})}_{\tilde{\mathcal{F}}^{-1}}
    = 
    \sup\limits_{\phi \in H^1_0\qty(\Omega) \backslash \{ 0 \} }
    \dfrac{\qty|\langle \mathcal{F}\qty(\hat{u}),\phi \rangle|}{\norm{\phi}_{\tilde{\mathcal{F}}}}
\end{align}
また，\(\bar{v} \coloneqq \mathcal{B}\hat{u}\)とし
\begin{align}
  \norm{\mathcal{F}\qty(\hat{u})}_{\tilde{\mathcal{F}}^{-1}}
  =&
  \sup\limits_{\phi \in H^1_0\qty(\Omega) \backslash \{ 0 \} }
  \dfrac{\qty|\langle \mathcal{F}\qty(\hat{u}),\phi \rangle|}{\norm{\phi}_{\tilde{\mathcal{F}}}}\\
  =&
  \sup\limits_{\phi \in H^1_0\qty(\Omega) \backslash \{ 0 \} }
  \dfrac{\qty|(\nabla \hat{u},\nabla \phi)_{L^2}-(f(\hat{u}),\phi)_{L^2}+\frac{\delta}{\epsilon^2}(\mathcal{B}\hat{u},\phi)_{L^2}|}{\norm{\phi}_{\tilde{\mathcal{F}}}}\\
  =&
  \sup\limits_{\phi \in H^1_0\qty(\Omega) \backslash \{ 0 \} }
  \dfrac{|(\nabla \hat{u},\nabla \phi)_{L^2}-(f(\hat{u}),\phi)_{L^2}
  +\frac{\delta}{\epsilon^2}(\bar{v},\phi)_{L^2}+(\frac{\delta}{\epsilon^2}\hat{v},\phi)_{L^2}-(\frac{\delta}{\epsilon^2}\hat{v},\phi)_{L^2}|}
  {\norm{\phi}_{\tilde{\mathcal{F}}}}\\
  \leq&
  \sup\limits_{\phi \in H^1_0\qty(\Omega) \backslash \{ 0 \} }
  \dfrac{|(\nabla \hat{u},\nabla \phi)_{L^2}-(f(\hat{u}),\phi)_{L^2}+(\frac{\delta}{\epsilon^2}\hat{v},\phi)_{L^2}|
  +\frac{\delta}{\epsilon^2}|(\bar{v},\phi)_{L^2}-(\hat{v},\phi)_{L^2}|}
  {\norm{\phi}_{\tilde{\mathcal{F}}}}\\
  \leq&
  \sup\limits_{\phi \in H^1_0\qty(\Omega) \backslash \{ 0 \} }
  \dfrac{\norm{-\Delta \hat{u}-f(\hat{u})+\frac{\delta}{\epsilon^2}\hat{v}}_{L^2}\norm{\phi}_{L^2}}
  {\norm{\phi}_{\tilde{\mathcal{F}}}}
  +\qty|\frac{\delta}{\epsilon^2}|
  \sup\limits_{\phi \in H^1_0\qty(\Omega) \backslash \{ 0 \} }
  \dfrac{\norm{   \bar{v}-\hat{v}}_{L^2}\norm{\phi}_{L^2}}
  {\norm{\phi}_{\tilde{\mathcal{F}}}}\\
  \leq&
  C_{p\tilde{\mathcal{F}}}\qty(
  \sup\limits_{\phi \in H^1_0\qty(\Omega) \backslash \{ 0 \} }
  \dfrac{\norm{-\Delta \hat{u}-f(\hat{u})+\frac{\delta}{\epsilon^2}\hat{v}}_{L^2}\norm{\phi}_{\tilde{\mathcal{F}}}}
  {\norm{\phi}_{\tilde{\mathcal{F}}}})\\
  +&
  \qty|\frac{\delta}{\epsilon^2}|
  C_{p\tilde{\mathcal{F}}}\qty(
  \sup\limits_{\phi \in H^1_0\qty(\Omega) \backslash \{ 0 \} }
  \dfrac{\norm{   \bar{v}-\hat{v}}_{L^2}\norm{\phi}_{\tilde{\mathcal{F}}}}
  {\norm{\phi}_{\tilde{\mathcal{F}}}}
  )\\
  =&
  C_{p\tilde{\mathcal{F}}}
  \qty(
    \norm{-\Delta \hat{u}-f(\hat{u})+\frac{\delta}{\epsilon^2}\hat{v}}_{L^2}
    +\qty|\frac{\delta}{\epsilon^2}|
    \norm{\bar{v}-\hat{v}}_{L^2}
  )
\end{align}
と評価できる．また，\(C_{p\tilde{\mathcal{F}}}\)は\(\norm{v}_{L^2}\leq C_{p\tilde{\mathcal{F}}}\norm{v}_{\tilde{\mathcal{F}}}\)を満たす定数とする．
ここで新しく\(r\)内積を以下のように定義する．
\begin{align}
   (u,v)_r \coloneqq (\nabla u,\nabla v)_{L^2}+r(u,v)_{L^2}
\end{align}
\(r\)内積から誘導される\(r\)ノルムを用いると，
\begin{align}
   \norm{\mathcal{F}\qty(\hat{u})}_{\tilde{\mathcal{F}}^{-1}}
   \leq&
   C_{p\tilde{\mathcal{F}}}\qty(\norm{-\Delta \hat{u}-f(\hat{u})+\dfrac{\delta}{\epsilon^2}\hat{v}}_{L^2}+\qty|\dfrac{\delta}{\epsilon^2}|\norm{\bar{v}-\hat{v}}_{r})\\
   \leq&
   C_{p\tilde{\mathcal{F}}}\qty(\norm{-\Delta \hat{u}-f(\hat{u})+\dfrac{\delta}{\epsilon^2}\hat{v}}_{L^2}+\qty|\dfrac{\delta}{\epsilon^2}|C_r\norm{-\hat{u}-(\Delta-r)\hat{v}}_{L^2})\\
\end{align}
この時，
\begin{align}
   &\norm{\bar{v}-\hat{v}}_r\\
   \leq&C_r\norm{-\hat{u}-(\Delta-r)\hat{v}}_{L^2}
\end{align}
となる\(C_r\)が存在する．
% \begin{align}
%    &\qty|(\nabla \hat{u},\nabla \phi)_{L^2}-(f(\hat{u}),\phi)_{L^2}+\frac{\delta}{\epsilon^2}(\mathcal{B}\hat{u},\phi)_{L^2}|\\
%    =&\qty|(\nabla \hat{u},\nabla \phi)_{L^2}-(f(\hat{u}),\phi)_{L^2}+\frac{\delta}{\epsilon^2}(\hat{v},\phi)_{L^2}|
%    +\qty|\frac{\delta}{\epsilon^2}(\mathcal{B}\hat{u},\phi)_{L^2}-\frac{\delta}{\epsilon^2}(\hat{v},\phi)_{L^2}|\\
% \end{align}
% \begin{align}
%    &\qty|\frac{\delta}{\epsilon^2}(\mathcal{B}\hat{u},\phi)_{L^2}-\frac{\delta}{\epsilon^2}(\hat{v},\phi)_{L^2}|\\
%    =&\qty|\frac{\delta}{\epsilon^2}|\qty|(\mathcal{B}\hat{u},\phi)_{L^2}-(\hat{v},\phi)_{L^2}|\\
%    =&\qty|\frac{\delta}{\epsilon^2}|\qty|(\bar{v},\phi)_{L^2}-(\hat{v},\phi)_{L^2}|\\
% \end{align}
ここで，
\begin{align}
   (\Delta-r)\bar{v}&=-\hat{u}\\
   (\Delta-r)(\bar{v}-\hat{v})&=-\hat{u}-(\Delta-r)\hat{v}\\
   \bar{v}-\hat{v}&=(\Delta-r)^{-1}(-\hat{u}-(\Delta-r)\hat{v})
\end{align}
であり，
\begin{align}
   &\norm{\bar{v}-\hat{v}}_r\\
   \leq&C_r\norm{(\Delta-r)((\Delta-r)^{-1}(-\hat{u}-(\Delta-r)\hat{v}))}_{L^2}\\
   =&C_r\norm{-\hat{u}-(\Delta-r)\hat{v}}_{L^2}
\end{align}
である\(C_r\)を求めたい．
\(\norm{\hat{v}-P_{h,r}v}_r\)について式変形をすると
\begin{align}
   &\norm{\hat{v}-P_{h,r}v}_r^2\\
   =&\norm{(I-P_{h,r})\hat{v}}_r^2\\
   \leq&\norm{(I-P_{h})\hat{v}}_{H^1_0}+r\norm{(I-P_{h})\hat{v}}_{L^2}\\
   \leq&\norm{(I-P_{h})\hat{v}}_{H^1_0}+rC_{N}^2\norm{(I-P_{h})\hat{v}}_{H^1_0}\\
   =&(1+rC_N^2)\norm{(I-P_{h})\hat{v}}_{H^1_0}\\
   \leq&(1+rC_N^2)C_N^2\norm{\Delta\hat{v}}_{L^2}\\
   \leq&(1+rC_N^2)C_N^2\norm{(\Delta-r)\hat{v}}_{L^2}\\
\end{align}
このことから，
\begin{align}
   C_r=C_N\sqrt{(1+rC_N^2)}
\end{align}
となる．
よって，
\begin{align}
   \norm{\bar{v}-\hat{v}}_r
   \leq
   C_N\sqrt{(1+rC_N^2)}\norm{-\hat{u}-(\Delta-r)\hat{v}}_{L^2}
\end{align}
となる．
よって\(\norm{\mathcal{F}'(\hat{u})}_{\tilde{\mathcal{F}}^{-1}}\)は
\begin{align}
   \norm{\mathcal{F}'(\hat{u})}_{\tilde{\mathcal{F}}^{-1}}
   \leq&
   C_{p\tilde{\mathcal{F}}}\qty(\norm{-\Delta \hat{u}-f(\hat{u})+\dfrac{\delta}{\epsilon^2}\hat{v}}_{L^2}+\qty|\dfrac{\delta}{\epsilon^2}|\qty(C_N\sqrt{(1+rC_N^2)})\norm{-\hat{u}-(\Delta-r)\hat{v}}_{L^2})\\
   \eqqcolon & R
\end{align}
となる．\\
まず\(\norm{-\Delta \hat{u}-f(\hat{u})+\dfrac{\delta}{\epsilon^2}\hat{v}}_{L^2}\)を計算する．
ノルムの定義より，\(\norm{-\Delta \hat{u}-f(\hat{u})+\dfrac{\delta}{\epsilon^2}\hat{v}}_{L^2}\)を２乗すると，
\begin{align}
   &\norm{-\Delta \hat{u}-(\dfrac{1}{\epsilon^2}(-\alpha \hat{u} + \beta \hat{u}^2 -\gamma \hat{u}^3-\delta \hat{v}))}_{L^2}^2\\
   &=(-\Delta \hat{u}-(\dfrac{1}{\epsilon^2}(-\alpha \hat{u} + \beta \hat{u}^2 -\gamma \hat{u}^3-\delta \hat{v})),
      -\Delta \hat{u}-(\dfrac{1}{\epsilon^2}(-\alpha \hat{u} + \beta \hat{u}^2 -\gamma \hat{u}^3-\delta \hat{v})))_{L^2}
\end{align}
となる．ここからそれぞれの項ごとに計算すると，
\begin{align}
   =&(-\Delta \hat{u},-\Delta \hat{u})_{L^2}\\
   +&\qty(-\Delta \hat{u},-\qty(\dfrac{1}{\epsilon^2}\qty(-\alpha \hat{u} + \beta \hat{u}^2 -\gamma \hat{u}^3-\delta \hat{v})))_{L^2}\\
   +&\qty(-\qty(\frac{1}{\epsilon^2}\qty(-\alpha \hat{u} + \beta \hat{u}^2 -\gamma \hat{u}^3-\delta \hat{v})),-\Delta \hat{u})_{L^2}\\
   +&\qty(-\qty(\frac{1}{\epsilon^2}\qty(-\alpha \hat{u} + \beta \hat{u}^2 -\gamma \hat{u}^3-\delta \hat{v})),-\qty(\frac{1}{\epsilon^2}(-\alpha \hat{u} + \beta \hat{u}^2 -\gamma \hat{u}^3-\delta \hat{v})))_{L^2}\\
\end{align}
ここで，
\(\qty(-\Delta \hat{u},-\qty(\frac{1}{\epsilon^2}(-\alpha \hat{u} + \beta \hat{u}^2 -\gamma \hat{u}^3-\delta \hat{v})))_{L^2}
+\qty(-\qty(\frac{1}{\epsilon^2}\qty(-\alpha \hat{u} + \beta \hat{u}^2 -\gamma \hat{u}^3-\delta \hat{v})),-\Delta \hat{u})_{L^2}\)
について式変形をすると，
\begin{align}
 &\qty(-\Delta \hat{u},-\qty(\frac{1}{\epsilon^2}(-\alpha \hat{u} + \beta \hat{u}^2 -\gamma \hat{u}^3-\delta \hat{v})))_{L^2}
 +\qty(-\qty(\frac{1}{\epsilon^2}(-\alpha \hat{u} + \beta \hat{u}^2 -\gamma \hat{u}^3-\delta \hat{v})),-\Delta \hat{u})_{L^2}\\
 &=
 \dfrac{2}{\epsilon^2}\{(-\Delta \hat{u},\alpha \hat{u})_{L^2}
 +(-\Delta \hat{u},-\beta \hat{u}^2)_{L^2}
 +(-\Delta \hat{u},\gamma \hat{u}^3)_{L^2}
 +(-\Delta \hat{u},\delta \hat{v})_{L^2}\}\\
 &=
 \dfrac{2}{\epsilon^2}\{\alpha(-\Delta \hat{u},\hat{u})_{L^2}
 -\beta(-\Delta \hat{u}, \hat{u}^2)_{L^2}
 +\gamma(-\Delta \hat{u}, \hat{u}^3)_{L^2}
 +\delta(-\Delta \hat{u}, \hat{v})_{L^2}\}\\
\end{align}
と変形できる．また，
\(\qty(-\qty(\dfrac{1}{\epsilon^2}\qty(-\alpha \hat{u} + \beta \hat{u}^2 -\gamma \hat{u}^3-\delta \hat{v})),
   -\qty(\dfrac{1}{\epsilon^2}\qty(-\alpha \hat{u} + \beta \hat{u}^2 -\gamma \hat{u}^3-\delta \hat{v})))_{L^2}\)
について式変形をすると，
\begin{align}
  &\qty(-\qty(\dfrac{1}{\epsilon^2}\qty(-\alpha \hat{u} + \beta \hat{u}^2 -\gamma \hat{u}^3-\delta \hat{v})),
   -\qty(\dfrac{1}{\epsilon^2}\qty(-\alpha \hat{u} + \beta \hat{u}^2 -\gamma \hat{u}^3-\delta \hat{v})))_{L^2}\\
  =&
  \qty(-\dfrac{1}{\epsilon^2}(-\alpha \hat{u}),-\frac{1}{\epsilon^2}(-\alpha \hat{u} + \beta \hat{u}^2 -\gamma \hat{u}^3-\delta \hat{v}))_{L^2}\\
  +&
  \qty(-\dfrac{1}{\epsilon^2}(\beta \hat{u}^2),-\frac{1}{\epsilon^2}(-\alpha \hat{u} + \beta \hat{u}^2 -\gamma \hat{u}^3-\delta \hat{v}))_{L^2}\\
  +&
  \qty(-\dfrac{1}{\epsilon^2}(-\gamma \hat{u}^3),-\frac{1}{\epsilon^2}(-\alpha \hat{u} + \beta \hat{u}^2 -\gamma \hat{u}^3-\delta \hat{v}))_{L^2}\\
  +&
  \qty(-\dfrac{1}{\epsilon^2}(-\delta \hat{v}),-\frac{1}{\epsilon^2}(-\alpha \hat{u} + \beta \hat{u}^2 -\gamma \hat{u}^3-\delta \hat{v}))_{L^2}\\
\end{align}
と変形できる．
各項ごとに計算すると\\
１項
\begin{align}
  &\qty(-\frac{1}{\epsilon^2}(-\alpha \hat{u}),-\dfrac{1}{\epsilon^2}(-\alpha \hat{u}))_{L^2}
  =\dfrac{1}{\epsilon^4}\alpha^2 \int \hat{u}(x)^2 dx\\
  &\qty(-\frac{1}{\epsilon^2}(-\alpha \hat{u}),-\dfrac{1}{\epsilon^2}(\beta \hat{u}^2))_{L^2}
  =-\dfrac{1}{\epsilon^4}\alpha\beta \int \hat{u}(x)^3 dx\\
  &\qty(-\frac{1}{\epsilon^2}(-\alpha \hat{u}),-\dfrac{1}{\epsilon^2}(-\gamma \hat{u}^3))_{L^2}
  =\dfrac{1}{\epsilon^4}\alpha\gamma \int \hat{u}(x)^4 dx\\
  &\qty(-\frac{1}{\epsilon^2}(-\alpha \hat{u}),-\dfrac{1}{\epsilon^2}(-\delta \hat{v}))_{L^2}
  =\dfrac{1}{\epsilon^4}\alpha\delta \int \hat{u}(x)\hat{v}(x) dx\\
\end{align}
２項
\begin{align}
  &\qty(-\dfrac{1}{\epsilon^2}(\beta \hat{u}^2),-\dfrac{1}{\epsilon^2}(-\alpha \hat{u}))_{L^2}
  =-\dfrac{1}{\epsilon^4}\alpha\beta \int \hat{u}(x)^3 dx\\
  &\qty(-\dfrac{1}{\epsilon^2}(\beta \hat{u}^2),-\dfrac{1}{\epsilon^2}(\beta \hat{u}^2))_{L^2}
  =\dfrac{1}{\epsilon^4}\beta^2 \int \hat{u}(x)^4 dx\\
  &\qty(-\dfrac{1}{\epsilon^2}(\beta \hat{u}^2),-\dfrac{1}{\epsilon^2}(-\gamma \hat{u}^3))_{L^2}
  =-\dfrac{1}{\epsilon^4}\beta\gamma \int \hat{u}(x)^5 dx\\
  &\qty(-\dfrac{1}{\epsilon^2}(\beta \hat{u}^2),-\dfrac{1}{\epsilon^2}(-\delta \hat{v}))_{L^2}
  =-\dfrac{1}{\epsilon^4}\beta\delta \int \hat{u}(x)^2\hat{v}(x) dx\\
\end{align}
３項
\begin{align}
  &\qty(-\dfrac{1}{\epsilon^2}(-\gamma \hat{u}^3),-\dfrac{1}{\epsilon^2}(-\alpha \hat{u}))_{L^2}
  =\dfrac{1}{\epsilon^4}\alpha\gamma \int \hat{u}(x)^4 dx\\
  &\qty(-\dfrac{1}{\epsilon^2}(-\gamma \hat{u}^3),-\dfrac{1}{\epsilon^2}(\beta \hat{u}^2))_{L^2}
  =-\dfrac{1}{\epsilon^4}\beta\gamma \int \hat{u}(x)^5 dx\\
  &\qty(-\dfrac{1}{\epsilon^2}(-\gamma \hat{u}^3),-\dfrac{1}{\epsilon^2}(-\gamma \hat{u}^3))_{L^2}
  =\dfrac{1}{\epsilon^4}\gamma^2 \int \hat{u}(x)^6 dx\\
  &\qty(-\dfrac{1}{\epsilon^2}(-\gamma \hat{u}^3),-\dfrac{1}{\epsilon^2}(-\delta \hat{v}))_{L^2}
  =\dfrac{1}{\epsilon^4}\gamma\delta \int \hat{u}(x)^3\hat{v}(x) dx\\
\end{align}
４項
\begin{align}
  &\qty(-\dfrac{1}{\epsilon^2}(-\delta \hat{v}),-\dfrac{1}{\epsilon^2}(-\alpha \hat{u}))_{L^2}
  =\dfrac{1}{\epsilon^4}\alpha\delta \int \hat{u}(x)\hat{v}(x) dx\\
  &\qty(-\dfrac{1}{\epsilon^2}(-\delta \hat{v}),-\dfrac{1}{\epsilon^2}(\beta \hat{u}^2))_{L^2}
  =-\dfrac{1}{\epsilon^4}\beta\delta \int \hat{u}(x)^2\hat{v}(x) dx\\
  &\qty(-\dfrac{1}{\epsilon^2}(-\delta \hat{v}),-\dfrac{1}{\epsilon^2}(-\gamma \hat{u}^3))_{L^2}
  =\dfrac{1}{\epsilon^4}\gamma\delta \int \hat{u}(x)^3\hat{v}(x) dx\\
  &\qty(-\dfrac{1}{\epsilon^2}(-\delta \hat{v}),-\dfrac{1}{\epsilon^2}(-\delta \hat{v}))_{L^2}
  =\dfrac{1}{\epsilon^4}\delta^2 \int \hat{v}(x)^2 dx\\
\end{align}
となる．したがって上記の式変形をまとめると
\begin{align}
  &\norm{-\Delta \hat{u}-(\dfrac{1}{\epsilon^2}(-\alpha \hat{u} + \beta \hat{u}^2 -\gamma \hat{u}^3-\delta \hat{v}))}_{L^2}^2\\
  =&(-\Delta \hat{u},-\Delta \hat{u})_{L^2}\\%first%
  +&\dfrac{2}{\epsilon^2}\{\alpha(-\Delta \hat{u},\hat{u})_{L^2}%second%
  -\beta(-\Delta \hat{u}, \hat{u}^2)_{L^2}
  +\gamma(-\Delta \hat{u}, \hat{u}^3)_{L^2}
  +\delta(-\Delta \hat{u}, \hat{v})_{L^2}\}\\
  +&\dfrac{1}{\epsilon^4}\delta^2 \int \hat{v}(x)^2 dx+\dfrac{2}{\epsilon^4}\alpha\delta \int \hat{u}(x)\hat{v}(x) dx\\
  +&\dfrac{1}{\epsilon^4}\alpha^2 \int \hat{u}(x)^2 dx
  -\dfrac{1}{\epsilon^4}2\beta\delta \int \hat{u}(x)^2\hat{v}(x) dx\\
  +&\dfrac{1}{\epsilon^4}2\alpha\beta\int \hat{u}(x)^3 dx
  +\dfrac{1}{\epsilon^4}2\gamma\delta \int \hat{u}(x)^3\hat{v}(x) dx\\
  +&\dfrac{1}{\epsilon^4}(2\alpha\gamma+\beta^2) \int \hat{u}(x)^4 dx -\dfrac{1}{\epsilon^4}2\beta\gamma \int \hat{u}(x)^5 dx +\dfrac{1}{\epsilon^4}\gamma^2 \int \hat{u}(x)^6 dx\\
\end{align}
となる．
次に\(\norm{-\hat{u}-(\Delta-r)\hat{v}}_{L^2}\)を計算する．ノルムの定義より
\begin{align}
  &\norm{-\hat{u}-(\Delta - r)\hat{v}}^2_{L^2}\\
  =&(-\hat{u}-(\Delta - r)\hat{v},-\hat{u}-(\Delta - r)\hat{v})_{L^2}\\
  =&(\hat{u},\hat{u})_{L^2}+(-(\Delta - r)\hat{v},-\hat{u})_{L^2}+(-\hat{u},-(\Delta - r)\hat{v})_{L^2}+(-(\Delta - r)\hat{v},-(\Delta - r)\hat{v})_{L^2}
\end{align}
この時，\((\hat{u},\hat{u})_{L^2}\)は
\begin{align}
  (\hat{u},\hat{u})_{L^2}=\int \hat{u}(x)^2 dx
\end{align}
となる．また，
\((-(\Delta - r)\hat{v},\hat{u})_{L^2}
  +(\hat{u},-(\Delta - r)\hat{v})_{L^2}\)は，
  \begin{align}
    &(-(\Delta - r)\hat{v},-\hat{u})_{L^2}+(-\hat{u},-(\Delta - r)\hat{v})_{L^2}\\
    &=2\qty((-\Delta \hat{v},-\hat{u})_{L^2}+(r \hat{v},-\hat{u})_{L^2})
  \end{align}
となる．
そして，\((-(\Delta - r)\hat{v},-(\Delta - r)\hat{v})_{L^2}\)は
\begin{align}
  &(-(\Delta - r)\hat{v},-(\Delta - r)\hat{v})_{L^2}\\
  &=(-\Delta \hat{v},-\Delta \hat{v})_{L^2}
  +(-\Delta \hat{v},r \hat{v})_{L^2}
  +(r \hat{v},-\Delta \hat{v})_{L^2}
  +(r \hat{v},r \hat{v})_{L^2}\\
  &=(-\Delta \hat{v},-\Delta \hat{v})_{L^2}
  -2r(\Delta \hat{v}, \hat{v})_{L^2}
  +r^2( \hat{v}, \hat{v})_{L^2}
\end{align}
まとめると，
% \norm{\hat{u}-(\Delta+r)\hat{v}}_{L^2}^2=
% (\hat{u},\hat{u})_{L^2}
% +2\qty((-\Delta \hat{v},\hat{u})_{L^2}-r( \hat{v},\hat{u})_{L^2})
% +r^2 (\hat{v},\hat{v})_{L^2}+(-\Delta \hat{v},-\Delta \hat{v})_{L^2}
% +2r(-\Delta \hat{v},- \hat{v})_{L^2}
% +r^2( \hat{v}, \hat{v})_{L^2}
\begin{align}
  &\norm{-\hat{u}-(\Delta-r)\hat{v}}_{L^2}^2\\
  &=
  (\hat{u},\hat{u})_{L^2}\\
  &+2\qty((-\Delta \hat{v},-\hat{u})_{L^2}+(r \hat{v},-\hat{u})_{L^2})\\
  &+(-\Delta \hat{v},-\Delta \hat{v})_{L^2}
  -2r(\Delta \hat{v}, \hat{v})_{L^2}
  +r^2( \hat{v}, \hat{v})_{L^2}
\end{align}
である．
よって，\(R\)は，\(C_{p\tilde{\mathcal{F}}}\)と
\begin{align}
  &(-\Delta \hat{u},-\Delta \hat{u})_{L^2}\\%first%
  +&\dfrac{2}{\epsilon^2}\{\alpha(-\Delta \hat{u},\hat{u})_{L^2}%second%
  -\beta(-\Delta \hat{u}, \hat{u}^2)_{L^2}
  +\gamma(-\Delta \hat{u}, \hat{u}^3)_{L^2}
  +\delta(-\Delta \hat{u}, \hat{v})_{L^2}\}\\
  +&\dfrac{1}{\epsilon^4}\delta^2 \int \hat{v}(x)^2 dx+\dfrac{2}{\epsilon^4}\alpha\delta \int \hat{u}(x)\hat{v}(x) dx\\
  +&\dfrac{1}{\epsilon^4}\alpha^2 \int \hat{u}(x)^2 dx
  -\dfrac{1}{\epsilon^4}2\beta\delta \int \hat{u}(x)^2\hat{v}(x) dx\\
  +&\dfrac{1}{\epsilon^4}2\alpha\beta\int \hat{u}(x)^3 dx
  +\dfrac{1}{\epsilon^4}2\gamma\delta \int \hat{u}(x)^3\hat{v}(x) dx\\
  +&\dfrac{1}{\epsilon^4}(2\alpha\gamma+\beta^2) \int \hat{u}(x)^4 dx 
  -\dfrac{1}{\epsilon^4}2\beta\gamma \int \hat{u}(x)^5 dx 
  +\dfrac{1}{\epsilon^4}\gamma^2 \int \hat{u}(x)^6 dx
\end{align}
の平方根と
\begin{align}
  C_r
  \qty((\hat{u},\hat{u})_{L^2}
  +2(-\Delta \hat{v},-\hat{u})_{L^2}+2(r \hat{v},-\hat{u})_{L^2}
  +(-\Delta \hat{v},-\Delta \hat{v})_{L^2}
  -2r(\Delta \hat{v}, \hat{v})_{L^2}
  +r^2( \hat{v}, \hat{v})_{L^2})
\end{align}
の平方根の和で求められる．


\section{\(G\)の評価}
%%%%%%%%%%%%%%%%% [[[[Gの評価]]]] %%%%%%%%%%%%%%%%%%%%%%%%%%%%%%%%
\(
    \norm{\mathcal{F}'\qty[w]-\mathcal{F}'\qty[l]}_{L\qty(W,W^{-1})}
    \leq
    G\norm{w-l}_{\tilde{\mathcal{F}}}
\)
を評価する．
\begin{align}
    \norm{\mathcal{F}'\qty[w]-\mathcal{F}'\qty[l]}_{L\qty(W,W^{-1})}
    =
    \sup\limits_{\phi \in H^1_0(\Omega) \backslash \{ 0 \} }
    \dfrac
    {\norm{\qty(\mathcal{F}'\qty[w]-\mathcal{F}'\qty[l])\phi}_{{\tilde{\mathcal{F}}}^{-1}}}
    {\norm{\phi}_{{\tilde{\mathcal{F}}}}} 
\end{align}
共役空間のノルム
\(
    \norm{T}_{{\tilde{\mathcal{F}}}^{-1}}=
    \sup\limits_{\phi \in H^1_0(\Omega) \backslash \{ 0 \} }
    \dfrac{|( T,\phi )_{L^2}|}{\norm{\phi}_{\tilde{\mathcal{F}}}}
\)
より
\begin{align}
    \norm{\mathcal{F}'\qty[w]-\mathcal{F}'\qty[l]}_{L\qty(W,W^{-1})}
    =&
    \sup\limits_{\phi \in H^1_0(\Omega) \backslash \{ 0 \} }
    \frac
    {\norm{\qty(\mathcal{F}'\qty[w]-\mathcal{F}'\qty[l])\phi}_{{\tilde{\mathcal{F}}}^{-1}}}
    {\norm{\phi}_{{\tilde{\mathcal{F}}}}}\\ 
    =&
    \sup\limits_{\phi \in H^1_0(\Omega) \backslash \{ 0 \} }
    \frac
    {\sup\limits_{\psi \in H^1_0(\Omega) \backslash \{ 0 \} }
    \frac
    {|\langle \mathcal{F}'\qty[w]-\mathcal{F}'\qty[l],\psi \rangle|}
    {\norm{\psi}_{\tilde{\mathcal{F}}}}}
    {\norm{\phi}_{{\tilde{\mathcal{F}}}}}\\
    =&
    \sup\limits_{\phi \in H^1_0(\Omega) \backslash \{ 0 \} }
    \sup\limits_{\psi \in H^1_0(\Omega) \backslash \{ 0 \} }
    \frac
    {|\langle \qty(\mathcal{F}'\qty[w]-\mathcal{F}'\qty[l])\phi,\psi \rangle|}
    {\norm{\psi}_{\tilde{\mathcal{F}}}\norm{\phi}_{\tilde{\mathcal{F}}}}
\end{align}
ここで
\begin{align}
    \mathcal{F}'\qty[w]-\mathcal{F}'\qty[l]
    =\mathcal{A} - \mathcal{N}'\qty[w]-\qty(\mathcal{A} - \mathcal{N}'\qty[l])
    =-\mathcal{N}'\qty[w]+\mathcal{N}'\qty[l]
\end{align}
であるので，
\begin{align}
    \norm{\mathcal{F}'\qty[w]-\mathcal{F}'\qty[l]}_{L\qty(W,W^{-1})}
    =
    \sup\limits_{\phi \in H^1_0(\Omega) \backslash \{ 0 \} }
    \sup\limits_{\psi \in H^1_0(\Omega) \backslash \{ 0 \} }
    \frac
    {|\langle -\qty(\mathcal{N}'\qty[w]+\mathcal{N}'\qty[l])\phi,\psi \rangle|}
    {\norm{\psi}_{\tilde{\mathcal{F}}}\norm{\phi}_{\tilde{\mathcal{F}}}}
\end{align}
と置き換えられる．また
\begin{align}
  % \mathcal{N}'\qty[w]=&\frac{1}{\epsilon^2}\qty(f^\prime[w]-\delta \mathcal{B})\\
  % \mathcal{N}'\qty[l]=&\frac{1}{\epsilon^2}\qty(f^\prime[l]-\delta \mathcal{B})
  &\langle \mathcal{N}'\qty[w]\phi,\psi \rangle =\qty(\frac{1}{\epsilon^2}\qty(f^\prime[w]-\delta \mathcal{B})\phi,\psi)_{L^2}\\
  &\langle \mathcal{N}'\qty[l]\phi,\psi \rangle =\qty(\frac{1}{\epsilon^2}\qty(f^\prime[l]-\delta \mathcal{B})\phi,\psi)_{L^2} 
\end{align}
より，
\begin{align}
  &\norm{\mathcal{F}'\qty[w]-\mathcal{F}'\qty[l]}_{L\qty(W,W^{-1})}\\
  &=
  \sup\limits_{\phi \in H^1_0(\Omega) \backslash \{ 0 \} }
  \sup\limits_{\psi \in H^1_0(\Omega) \backslash \{ 0 \} }
  \frac
  {\qty| \qty(\qty(-\frac{1}{\epsilon^2}\qty(f^\prime[w]-\delta \mathcal{B}) 
  + \frac{1}{\epsilon^2}\qty(f^\prime[l]-\delta \mathcal{B}))\phi,\psi)_{L^2} |}
  {\norm{\psi}_{\tilde{\mathcal{F}}}\norm{\phi}_{\tilde{\mathcal{F}}}}\\
  &=
  \sup\limits_{\phi \in H^1_0(\Omega) \backslash \{ 0 \} }
  \sup\limits_{\psi \in H^1_0(\Omega) \backslash \{ 0 \} }
  \frac
  {\qty| \qty(\qty(\frac{1}{\epsilon^2}\qty(f^\prime[l]-f^\prime[w]))\phi,\psi)_{L^2} |}
  {\norm{\psi}_{\tilde{\mathcal{F}}}\norm{\phi}_{\tilde{\mathcal{F}}}}\\
\end{align}
この時，
\begin{align}
  &f[l]=-\alpha l + \beta l^2 - \gamma l^3\\
  &f[w]=-\alpha w + \beta w^2 - \gamma w^3\\
\end{align}
とすると，
\begin{align}
  &f^\prime[l]=-\alpha  + \beta l - \gamma l^2\\
  &f^\prime[w]=-\alpha  + \beta w - \gamma w^2\\
\end{align}
であるので,
\begin{align}
  &\norm{\mathcal{F}'\qty[w]-\mathcal{F}'\qty[l]}_{L\qty(W,W^{-1})}\\
  &=
  \sup\limits_{\phi \in H^1_0(\Omega) \backslash \{ 0 \} }
  \sup\limits_{\psi \in H^1_0(\Omega) \backslash \{ 0 \} }
  \frac
  {\qty| (\qty(\frac{1}{\epsilon^2}\qty(f^\prime[l]-f^\prime[w]))\phi,\psi)_{L^2} |}
  {\norm{\psi}_{\tilde{\mathcal{F}}}\norm{\phi}_{\tilde{\mathcal{F}}}}\\
  &=
  \sup\limits_{\phi \in H^1_0(\Omega) \backslash \{ 0 \} }
  \sup\limits_{\psi \in H^1_0(\Omega) \backslash \{ 0 \} }
  \frac
  {\qty|( \qty(\frac{1}{\epsilon^2}
  \qty(\qty(-\alpha  + \beta l - \gamma l^2)-\qty(-\alpha  + \beta w - \gamma w^2)))\phi,\psi )_{L^2}|}
  {\norm{\psi}_{\tilde{\mathcal{F}}}\norm{\phi}_{\tilde{\mathcal{F}}}}\\
  &=
  \sup\limits_{\phi \in H^1_0(\Omega) \backslash \{ 0 \} }
  \sup\limits_{\psi \in H^1_0(\Omega) \backslash \{ 0 \} }
  \frac
  {\qty|( \qty(\frac{1}{\epsilon^2}
  \qty(\beta\qty(l-w)-\gamma\qty(l+w)\qty(l-w)))\phi,\psi )_{L^2}|}
  {\norm{\psi}_{\tilde{\mathcal{F}}}\norm{\phi}_{\tilde{\mathcal{F}}}}\\
  &\leq
  \sup\limits_{\phi \in H^1_0(\Omega) \backslash \{ 0 \} }
  \sup\limits_{\psi \in H^1_0(\Omega) \backslash \{ 0 \} }
  \frac
  {\frac{\beta}{\epsilon^2}\norm{l-w}_{L^3}
  \norm{\psi}_{L^3}\norm{\phi}_{L^3}}
  {\norm{\psi}_{\tilde{\mathcal{F}}}\norm{\phi}_{\tilde{\mathcal{F}}}}\\
  &+
  \sup\limits_{\phi \in H^1_0(\Omega) \backslash \{ 0 \} }
  \sup\limits_{\psi \in H^1_0(\Omega) \backslash \{ 0 \} }
  \frac
  {\frac{\gamma}{\epsilon^2}\norm{l-w}_{L^4}\norm{l+w}_{L^4}
  \norm{\psi}_{L^4}\norm{\phi}_{L^4}}
  {\norm{\psi}_{\tilde{\mathcal{F}}}\norm{\phi}_{\tilde{\mathcal{F}}}}\\
\end{align}
ここでソボレフの埋め込み定理より，
\begin{align}
  \norm{u}_{L^3}\leq C_{{\tilde{\mathcal{F}}},3} \norm{u}_{H^1_0}\leq C_{{\tilde{\mathcal{F}}},3}\norm{u}_{\tilde{\mathcal{F}}}\\
  \norm{u}_{L^4}\leq C_{{\tilde{\mathcal{F}}},4} \norm{u}_{H^1_0}\leq C_{{\tilde{\mathcal{F}}},4}\norm{u}_{\tilde{\mathcal{F}}}\\
\end{align}
が成り立つ．よって
\begin{align}
  &\norm{\mathcal{F}'\qty[w]-\mathcal{F}'\qty[l]}_{L\qty(W,W^{-1})}\\
  \leq&
  \sup\limits_{\phi \in H^1_0(\Omega) \backslash \{ 0 \} }
  \sup\limits_{\psi \in H^1_0(\Omega) \backslash \{ 0 \} }
  \frac
  {\frac{\beta}{\epsilon^2}\norm{l-w}_{L^3}\norm{\phi}_{L^3}\norm{\psi}_{L^3}}
  {\norm{\psi}_{\tilde{\mathcal{F}}}\norm{\phi}_{\tilde{\mathcal{F}}}}\\
  -&
  \sup\limits_{\phi \in H^1_0(\Omega) \backslash \{ 0 \} }
  \sup\limits_{\psi \in H^1_0(\Omega) \backslash \{ 0 \} }
  \frac
  {\frac{\gamma}{\epsilon^2}\norm{l-w}_{L^4}\norm{l+w}_{L^4}\norm{\phi}_{L^4}\norm{\psi}_{L^4}}
  {\norm{\psi}_{\tilde{\mathcal{F}}}\norm{\phi}_{\tilde{\mathcal{F}}}}\\
  \leq&
  \sup\limits_{\phi \in H^1_0(\Omega) \backslash \{ 0 \} }
  \sup\limits_{\psi \in H^1_0(\Omega) \backslash \{ 0 \} }
  \frac
  {\frac{\beta}{\epsilon^2}C^3_{{\tilde{\mathcal{F}}},3}\norm{l-w}_{\tilde{\mathcal{F}}}\norm{\phi}_{\tilde{\mathcal{F}}}\norm{\psi}_{\tilde{\mathcal{F}}}}
  {\norm{\psi}_{\tilde{\mathcal{F}}}\norm{\phi}_{\tilde{\mathcal{F}}}}\\
  -&
  \sup\limits_{\phi \in H^1_0(\Omega) \backslash \{ 0 \} }
  \sup\limits_{\psi \in H^1_0(\Omega) \backslash \{ 0 \} }
  \frac
  {\frac{\gamma}{\epsilon^2}C^4_{{\tilde{\mathcal{F}}},4}\norm{l-w}_{\tilde{\mathcal{F}}}\norm{l+w}_{\tilde{\mathcal{F}}}\norm{\phi}_{\tilde{\mathcal{F}}}\norm{\psi}_{\tilde{\mathcal{F}}}}
  {\norm{\psi}_{\tilde{\mathcal{F}}}\norm{\phi}_{\tilde{\mathcal{F}}}}\\
  =&
  \frac{\beta}{\epsilon^2}
  \qty(C^3_{{\tilde{\mathcal{F}}},3}\norm{l-w}_{\tilde{\mathcal{F}}})+\frac{\gamma}{\epsilon^2}\qty(C^4_{{\tilde{\mathcal{F}}},4}\norm{l-w}_{\tilde{\mathcal{F}}}\norm{l+w}_{\tilde{\mathcal{F}}})\\
  \leq&
  \frac{\beta}{\epsilon^2}
  \qty(C^3_{{\tilde{\mathcal{F}}},3}\norm{l-w}_{\tilde{\mathcal{F}}})+\frac{\gamma}{\epsilon^2}\qty(C^4_{{\tilde{\mathcal{F}}},4}(2\norm{\hat{u}}_{\tilde{\mathcal{F}}}+4a)\norm{l-w}_{\tilde{\mathcal{F}}})\\
\end{align}
したがって
%\(\frac{\beta}{\epsilon^2}
%  \qty(C^3_{{\tilde{\mathcal{F}}},3}\norm{l-w}_{\tilde{\mathcal{F}}})+\frac{\gamma}{\epsilon^2}\qty(C^4_{{\tilde{\mathcal{F}}},4}(2\norm{\hat{u}}_{\tilde{\mathcal{F}}}+4a)\norm{l-w}_{\tilde{\mathcal{F}}})\)となる．
\(G\coloneqq \frac{\beta}{\epsilon^2}
  C^3_{\tilde{\mathcal{F}},3}+\frac{\gamma}{\epsilon^2}\qty(C^4_{\tilde{\mathcal{F}},4}(2\norm{\hat{u}}_{\tilde{\mathcal{F}}}+4a))\)となる．

ただし，\(\norm{l\pm w}_{\tilde{\mathcal{F}}}\)は\(l,w\in \bar{B}\qty(\hat{u},2a)\)から評価できる．

% 提案手法を書いてね
